\clearpage
\thispagestyle{empty}\SourceLeaf{283}
\pdfbookmark[0]{Anzeigen eigner Abhandlungen}{anzeigen-title}
\begin{center}
\vspace*{2in}
{\Large ANZEIGEN\par}
\vspace{0.5in}
EIGNER\par
\vspace{0.5in}
{\Large ABHANDLUNGEN.\par}
\end{center}
\clearpage
\SourceLeaf{285}
\fancyhead[L]{\small Anzeigen eigner Abhandlungen}
\pdfbookmark[1]{Göttingische gelehrte Anzeigen. 1813 April 5.}{anzeige-1813-april5}
\begin{center}
\rule{0.85\textwidth}{0.3pt}\par
Göttingische gelehrte Anzeigen. 1813 April 5.\par
\rule{0.85\textwidth}{0.3pt}
\end{center}
\medskip
Am 18. März wurde der königl. Societät von dem Prof. \textsc{Gauss} eine Vorlesung übergeben, mit der Überschrift:
\begin{center}
\textit{Theoria attractionis corporum sphaeroidicorum ellipticorum homogeneorum, methodo nova tractata.}
\end{center}
Bekanntlich haben sich seit \textsc{Newton}'s Zeiten die ersten Geometer mit dieser berühmten Aufgabe gleichsam wetteifernd beschäftigt. \textsc{Newton} that den ersten Schritt, indem er die Anziehung eines Punktes in der Axe eines Umdrehungs-Sphäroids bestimmen lehrte, so wie ausserdem das einfache Verhältniss zwischen den Anziehungen aller Punkte, die im Innern des Sphäroids in einem und demselben Diameter liegen. \textsc{Maclaurin} glückte es hiernächst, in seiner berühmten Preisschrift über die Ebbe und Fluth die Anziehung aller Punkte auf der Oberfläche des Sphäroids zu bestimmen, auf welche sich auch, vermöge des \textsc{Newton}schen Lehrsatzes, die Anziehungen im Innern reduciren liessen, so dass also nur noch die Anziehung der äussern Punkte fehlte, deren Bestimmung freilich den schwierigsten Theil der Aufgabe ausmachte. Auch hierin that \textsc{Maclaurin} schon einen Schritt: er bestimmte die Anziehung der Punkte in der Verlängerung des Aequators und der Axe. \textsc{Maclaurin}'s Entdeckungen wurden als Meisterstücke der Synthesis allgemein bewundert, und eine Zeitlang als Beweise angesehen, dass es Fälle gebe, wo die synthetische Methode einen entschiedenen Vorzug vor
\SourceLeaf{286}der analytischen habe. \textsc{Lagrange} setzte letztere wieder in ihre Rechte ein, indem er ihr eine Aufgabe unterwürfig machte, welche nur der Synthesis zugänglich geschienen hatte, und mit der ihm eignen Gewandtheit alle Entdeckungen \textsc{Maclaurin}'s auf analytischem Wege zu finden lehrte. Obgleich dadurch in der Sache selbst kein neuer Fortschritt gemacht war, so musste dies doch als eine höchst wichtige Vorbereitung der spätern Arbeiten angesehen werden. \textsc{Legendre} war es, dem es gelang, die Theorie der Anziehung der Umdrehungs-Sphäroide zu vollenden, indem er den schönen Lehrsatz fand, und bewies, dass die Anziehung eines äussern Punktes von einem Sphäroide dieselbe Richtung hat, wie die Anziehung desselben Punktes von einem zweiten Sphäroide, dessen Oberfläche durch diesen Punkt geht, wenn die beiden erzeugenden Ellipsen einerlei Brennpunkt haben, und dass die erstere Anziehung sich zur andern verhält, wie die Masse des erstern Sphäroids zur Masse des andern.

Alles dieses bezieht sich auf die Sphäroide, welche durch Umdrehung einer halben Ellipse um die eine oder die andere Axe erzeugt sind. Allein jetzt blieb noch die weit schwerere Aufgabe zurück, die Anziehung eines Ellipsoids zu bestimmen, bei welchem auch der Aequator elliptisch ist, oder eines Körpers, von welchem jeder Schnitt mit einer Ebene eine Ellipse gibt. Die Bestimmung der Anziehung für Punkte in der Richtung der drei Hauptaxen hatte schon \textsc{Maclaurin} angedeutet, und \textsc{d'Alembert} und \textsc{Lagrange} hatten dafür analytische Beweise gegeben. \textsc{Legendre} hatte ferner aus Induction die allgemeine Gültigkeit seines vorhin angeführten schönen Theorems geahnt, ohne doch einen strengen Beweis finden zu können. \textsc{Laplace} war es vorbehalten, diese Lücke auszufüllen, und die Auflösung der Aufgabe in ihrer ganzen Allgemeinheit zu vollenden (1782). Hiermit, könnte man glauben, sei nun die Untersuchung als geschlossen anzusehen. Allein schon der Umstand, dass mehrere Geometer seit der Zeit sich wieder von neuem mit demselben Gegenstande beschäftigt haben, zeigt, dass noch viel zu wünschen übrig blieb. Das Erste und Wesentlichste bei einer Aufgabe ist immer, dass sie überhaupt nur aufgelöst werde. Allein zu einem und demselben Ziele führen oft mehrere Wege. Nicht selten kommt man zum ersten Male auf einem langen dornigen Umwege zum Ziele; der kürzeste, der wahre echte Weg wird erst viel später entdeckt. Die \textsc{Laplace}'sche Auflösung ist ein schönes Document der feinsten analytischen Kunst: allein der Weg, auf welchem er dazu gelangt. ist lang und beschwerlich, und gewiss ist die Anzahl der Geometer und
\SourceLeaf{287}Astronomen, die ihm darauf gefolgt sind, nur klein. Auch der Gebrauch der unendlichen Reihen, deren Convergenz nicht bewiesen ist, thut der Klarheit und Bündigkeit des Beweises einigen Eintrag. \textsc{Legendre} hat zwar 1788 eine andere Auflösung gegeben, von welcher indess fast dasselbe gilt, was wir gegen die von \textsc{Laplace} erinnert haben. Ein competenter Richter, \textsc{Lagrange}, fällt über die Auflösungen jener beiden grossen Analysten folgendes Urtheil (in den \textit{Nouv. Mém. de l'Acad. de Berlin} 1793): ‘\textit{On ne peut regarder leurs solutions que comme des chefs-d'oeuvre d'analyse, mais on peut désirer encore une solution plus directe et plus simple: et les progrès continuels de l'analyse donnent lieu de l'espérer.}’ Seitdem haben noch \textsc{Biot} und \textsc{Plana} jene beiden Beweise zu vervollkommnen und zu vereinfachen gesucht. Indessen obgleich diese Arbeiten schätzbar sind, muss man doch noch immer diese Auflösungen zu den verwickeltsten und subtilsten Anwendungen der Analyse rechnen.

Der Verfasser der gegenwärtigen Abhandlung, welcher seit lange schon die Überzeugung hatte, dass die \textit{echte} Auflösungsmethode jener berühmten Aufgabe erst noch gefunden werden müsse, wurde vor einem halben Jahre veranlasst, sich mit derselben näher zu beschäftigen, und indem er einen von den vorigen ganz abweichenden Weg nahm, hatte er das Vergnügen, auf eine so überraschend kurze und einfache Auflösung zu kommen, dass das Wesentliche davon sich auf zwei Seiten bringen liess. Freilich hat er sie hier nicht ganz so kurz vorgetragen. Theils wünschte er sie auch weniger geübten Lesern verständlich zu machen (denen diese für die Gestalt der Erde so interessanten Untersuchungen bisher ganz unzugänglich waren), und dass sich die neue Auflösung dazu vollkommen qualificire, davon hat er bereits mehrere Beweise. Theils schien es der Mühe werth, die Gründe, worauf sie beruht, und die auch bei andern Gelegenheiten oft mit Vortheil anzuwenden sein werden, etwas ausführlicher zu entwickeln, als für den nächsten Zweck erforderlich gewesen wäre.

Wir wollen jetzt hier noch die Hauptmomente der ganzen Auflösung in möglichster Kürze darstellen, doch für Kenner vollkommen hinlänglich. Wir müssen hier Verzicht darauf leisten, auch solchen Lesern ganz verständlich zu werden, die mit Untersuchungen dieser Art noch nicht vertraut sind; diese müssen wir auf die ausführliche Abhandlung selbst verweisen, welche schon gedruckt ist, und in kurzem in dem zweiten Bande der \textit{Commentationes recentiores} der Societät erscheinen wird.

\SourceLeaf{288}Der Verf. fängt damit an, sechs verschiedene allgemeine Lehrsätze zu begründen, vermittelst deren dreifache, durch einen körperlichen Raum auszudehnende, Integrale auf zweifache, nur über die Oberfläche des Körpers auszudehnende, Integrale reducirt werden. Wir geben hier von diesen Lehrsätzen nur drei, da die andern zur gegenwärtigen Untersuchung nicht nothwendig sind.

Es sei $\mathrm{d}s$ ein Element der Oberfläche eines Körpers von beliebiger Gestalt; $PQ,PM,PX,PY,PZ$, gerade Linien, von einem Punkte $P$ dieses Elements gezogen, senkrecht auf die Oberfläche und nach aussen zu, nach dem angezogenen Punkte $M$, parallel mit den drei Axen der Coordinaten. Es sei ferner $r$ der Abstand des Punktes $M$ von $P$; $MQ$ der Winkel zwischen $PM$ und $PQ$; $MX$ der Winkel zwischen $PM$ und $PX$; $QX$ der Winkel zwischen $PQ$ und $PX$. Endlich bezeichne $\pi$ das Verhältniss des Kreisumfanges zum Durchmesser, $X$ die Anziehung, welche der ganze Körper auf den Punkt $M$ parallel mit den Coordinaten $x$ ausübt. Man hat dann
\[
\text{I.}\qquad \int\frac{\mathrm{d}s.\cos MQ}{rr}=0\text{ oder }=-4\pi
\]
je nachdem $M$ ausserhalb oder innerhalb des Körpers fällt,
\[
\begin{aligned}
\text{II.}\qquad &\int\frac{\mathrm{d}s.\cos QX}{r}=X\\
\text{III.}\qquad &\int\frac{\mathrm{d}s.\cos MQ.\cos MX}{r}=-X
\end{aligned}
\]
wo die Integrale über die ganze Oberfläche des Körpers auszudehnen sind. Die Beweise dieser Lehrsätze unterdrücken wir hier, und bemerken nur, dass die zwei ersten sich auf Zerlegung des Körpers in Kegelelemente, die ihre Spitze in $M$ haben, gründen, der dritte hingegen auf Zerlegung des Körpers in prismatische Elemente, parallel mit der Axe der Coordinaten $x$.

Für die Oberfläche eines Ellipsoids, dessen drei halbe Axen $A,B,C$ sind, hat man zwischen den Coordinaten $x,y,z$ die Gleichung
\[
\frac{xx}{AA}+\frac{yy}{BB}+\frac{zz}{CC}=1
\]
Ferner wird $\cos QX=\frac{x}{AA\rho}$, wenn man Kürze halber setzt
\[
\BareRoot{(\frac{xx}{A^4}+\frac{yy}{B^4}+\frac{zz}{C^4})}=\rho
\]
Bedeuten $a,b,c$ die Coordinaten des Punkts $M$, so hat man
\SourceLeaf{289}
\[
\begin{aligned}
r&=\BareRoot{[(a-x)^2+(b-y)^2+(c-z)^2]}\\
\cos MX&=\frac{a-x}{r}\\
\cos MQ&=\frac{1}{\rho r}(\frac{(a-x)x}{AA}+\frac{(b-y)y}{BB}+\frac{(c-z)z}{CC})
\end{aligned}
\]
Es werden jetzt zwei neue veränderliche Grössen $p,q$ eingeführt, von denen $x,y,z$ so abhangen, dass
\[
\begin{aligned}
x&=A\cos p\\
y&=B\sin p\cos q\\
z&=C\sin p\sin q
\end{aligned}
\]
Um also die ganze Oberfläche des Ellipsoids zu umfassen, muss man $p$ von $0$ bis $180^0$, $q$ von $0$ bis $360^0$ ausdehnen. Man setze endlich noch $X=ABC\xi$. Aus bekannten Gründen ergibt sich $\mathrm{d}s=\mathrm{d}p.\mathrm{d}q.ABC\rho\sin p$. Obige drei Theoreme erhalten hiedurch folgende Gestalt, wenn man Kürze halber
\[
\frac{(a-x)x}{AA}+\frac{(b-y)y}{BB}+\frac{(c-z)z}{CC}=\psi
\]
setzt,
\[
\begin{aligned}
[1]\qquad &\int\int\frac{\mathrm{d}p.\mathrm{d}q.\sin p.\psi}{r^3}=0\text{ oder }=-\frac{4\pi}{ABC}\\
[2]\qquad &\int\int\frac{\mathrm{d}p.\mathrm{d}q.\cos p.\sin p}{r}=A\xi\\
[3]\qquad &\int\int\frac{\mathrm{d}p.\mathrm{d}q.\sin p.\psi(a-x)}{r^3}=-\xi
\end{aligned}
\]
Man betrachte nun $A,B,C$ als bestimmte besondere Werthe dreier veränderlicher Grössen $\alpha,\beta,\gamma$, die aber so verbunden sind, dass $\alpha\alpha-\beta\beta$, $\alpha\alpha-\gamma\gamma$ constant bleiben. Die Formel [1] führt leicht zu dem Schluss, dass, für ein unendlich wachsendes $\alpha$, $\xi$ unendlich abnimmt. Differentiirt man [2] in Beziehung auf die veränderlichen Grössen $\alpha,\beta,\gamma$, und bedient sich dabei des Variationszeichens $\delta$, so kommt
\[
\begin{aligned}
\alpha\delta\xi+\xi\delta\alpha&=-\int\int\frac{\mathrm{d}p.\mathrm{d}q.\cos p.\sin p.\delta r}{rr}\\
&=\delta\alpha.\int\int\frac{\mathrm{d}p.\mathrm{d}q.\sin p.x\psi}{r^3}
\end{aligned}
\]
\SourceLeaf{290}oder wenn man hier statt $\xi$ seinen Werth aus [3] setzt,
\[
\alpha\delta\xi=\delta\alpha.\int\int\frac{\mathrm{d}p.\mathrm{d}q.\sin p.a\psi}{r^3}
\]

Dies mit [1] verglichen, gibt
\[
[4]\qquad\delta\xi=0
\]
wenn der Punkt $M$ ausserhalb des Ellipsoids,
\[
[5]\qquad\delta\xi=-\frac{4\pi a\delta\alpha}{\alpha\alpha\beta\gamma}
\]
wenn $M$ innerhalb liegt.

Aus [4] folgt, dass $\xi$ constant, oder die Anziehung der Masse proportional ist, für alle Sphäroide, deren Hauptschnitte Ellipsen von einerlei Brennpunkten sind, so lange $M$ nicht innerhalb fällt. Die Bestimmung der Anziehung eines Sphäroids auf einen äussern Punkt reducirt sich also auf die Bestimmung der Anziehung eines andern Sphäroids, das aus denselben Brennpunkten beschrieben durch den angezogenen Punkt geht. Um diese zu bestimmen, werde der andere Fall betrachtet, wo der angezogene Punkt innerhalb liegt. Durch die Substitution von
\[
\begin{aligned}
\beta\beta&=\alpha\alpha-AA+BB\\
\gamma\gamma&=\alpha\alpha-AA+CC
\end{aligned}
\]
in der Gleichung [5] wird diese, wenn man zugleich $\frac{A}{\alpha}=t$ setzt, und statt des Zeichens $\delta$ wieder das gewöhnliche $d$ schreibt
\[
\xi=\frac{4\pi a}{A^3}\int\frac{tt\mathrm{d}t}{\BareRoot{[(1-(1-\frac{BB}{AA})tt)(1-(1-\frac{CC}{AA})tt)]}}
\]
wo das Integral so bestimmt werden muss, dass es für $t=0$ verschwindet, und dann, für das \textit{bestimmte} Sphäroid, bis $t=1$ auszudehnen ist. Man hat also, in demselben Sinne,
\[
[6]\qquad X=\frac{4\pi aBC}{AA}\int\frac{tt\mathrm{d}t}{\BareRoot{[(1-(1-\frac{BB}{AA})tt)(1-(1-\frac{CC}{AA})tt)]}}
\]
Diese Formel gibt die Anziehung für alle Punkte, die nicht ausserhalb liegen, und da sie bis zur Oberfläche selbst gültig sein muss, und die Anziehung äusse\SourceLeaf{291}rer Punkte bereits auf die Anziehung der Punkte auf der Oberfläche zurückgeführt war, so ist dadurch die Aufgabe vollständig aufgelöst. (Es braucht kaum erinnert zu werden, dass die Anziehungen parallel mit den beiden andern Hauptaxen sich schlechthin durch Umtauschung von $A,a$ gegen $B,b$ oder gegen $C,c$ ergibt.)

Die Gleichung [6] lehrt ferner, dass für einen innern Punkt die Anziehung aller Sphäroide, die einander ähnlich sind und ähnlich liegen, identisch ist. Denkt man sich also ein solches Sphäroid in Schichten getheilt, die durch ähnliche ellipsoidische Flächen begrenzt sind, so ist klar, dass alle ausserhalb des angezogenen Punkts liegenden Schichten gar nichts zur Anziehung beitragen, und bloss die Anziehung des sphäroidischen Kerns übrig bleibt, dessen Oberfläche durch den angezogenen Punkt geht.

Zum Schluss erwähnt der Verf. noch der neuesten Arbeit über denselben Gegenstand von Hrn. \textsc{Ivory} in den \textit{Philos. Transact.} 1809, welche er, aufmerksam gemacht durch den Hrn. Grafen \textsc{Laplace}, erst kennen lernte, als seine eigene Abhandlung schon ganz vollendet war. Durch eine sehr glückliche Idee hat Hr. \textsc{Ivory} die Anziehung eines äussern Punktes auf die Anziehung eines innern zurückgeführt. Allein die Art, wie er die Anziehung innerer Punkte selbst bestimmt, ist zwar voll Scharfsinn und Kunst, aber zum Theil, eben so wie \textsc{Laplace}'s Auflösung für äussere Punkte, auf die Betrachtung unendlicher, nicht überall convergirender, Reihen gegründet, und weit von der Einfachheit entfernt, die gewünscht werden konnte, so dass die \textsc{Ivory}'sche Auflösung des Problems, als ein Ganzes betrachtet, im Grunde nicht viel weniger künstlich und verwickelt ist, als die Auflösungen von \textsc{Laplace} und \textsc{Legendre}. Übrigens beruhet jene und diejenige, von welcher hier Bericht erstattet ist, auf ganz verschiedenen Gründen, und beide haben gar nichts gemein, als den Gebrauch der zwei veränderlichen Grössen, welche oben mit $p,q$ bezeichnet sind.

\bigskip
\begin{center}
[\textit{Handschriftliche Bemerkung.}]
\end{center}
\medskip
Durch eine der hier vorgetragenen ähnliche Methode bestimmt man auch die Function $V$, d. i. die Summe aller Theilchen eines Körpers, jedes mit seinem Abstande vom angezogenen Punkt dividirt.
